% META: {"id": "TSPE-CONTIN-EX-005", "chapitre": "TSPE-CONTINUITE", "type_objet": "exercice", "capacites": ["TSPE-CONTIN-C1"], "capacites_codes": ["C1"], "methodes": ["M1"], "parcours": 3, "competences": ["raisonner", "chercher"], "duree_min": 15, "mode_creation": "ex_nihilo", "sources_inspiration": [], "fichier_tex": "chapitres/TSPE-CONTINUITE/exercices/TSPE-CONTIN-EX-005.tex", "corrige_tex": "chapitres/TSPE-CONTINUITE/corriges/TSPE-CONTIN-CO-005.tex", "status": "approved"}
% BEGIN-VERIFY
% from sympy import symbols, diff, nsolve
% x = symbols('x')
% f = x**5 + x - 1
% fp = diff(f, x)
% assert fp == 5*x**4 + 1
% assert f.subs(x, 0) == -1
% assert f.subs(x, 1) == 1
% c = nsolve(f, x, 0.7)
% assert 0.75 < float(c) < 0.76
% END-VERIFY
\begin{exercice}{TSPE-CONTIN-EX-005}{3}{15}
On considere la fonction $f(x) = x^5 + x - 1$ sur $\mathbb{R}$.
\begin{enumerate}
  \item Etudier le sens de variation de $f$ sur $\mathbb{R}$ (on justifiera que $f'(x) > 0$ pour tout $x$).
  \item Montrer que l'equation $f(x) = 0$ admet une unique solution reelle $\alpha$.
  \item Justifier que $\alpha \in [0, 1]$.
  \item Ecrire (sans l'executer) un algorithme Python de dichotomie qui renvoie un encadrement de $\alpha$ a $10^{-3}$ pres.
\end{enumerate}
\end{exercice}
